\section{Walsh-layer permutations}\label{sec:walsh}

This section and the next replace the Bluestein middle transform and chunk-exchange layout of~\cite{OAI26} by a direct
complex tensor FFT. The key new tool is an exact procedure that permutes the \emph{positions} of a binary record
address using simultaneous butterfly layers, per-record diagonal phases, and dyadic scalings, together with elementary
two-digit permutations at fewer than $K$ leftover positions outside the chunks (Section~\ref{sec:walsh-perm}). Three
full layers apply $H_0=\frac12\bigl(\begin{smallmatrix}1&1\\1&-1\end{smallmatrix}\bigr)$ exactly to any subset of the
selected bits (Section~\ref{sec:walsh-select}), two such selective layers around a sign diagonal execute a batch of
disjoint controlled flips (Section~\ref{sec:walsh-cnot}), and six stages of batched flips execute every positional
permutation (Section~\ref{sec:walsh-perm}).

Throughout this section and Section~\ref{sec:direct-fft}, stored arrays lie on the $q$-grid
$\Cg_q=2^{-q}(\Z+i\Z)$, which the bridge of Section~\ref{sec:bridge} instantiates at $q=P+2$. For arrays of records,
$\|\cdot\|_\infty$ denotes the largest complex modulus of a coefficient, and a \emph{disk array} is one with
$\|z\|_\infty\le1$. By~\eqref{eq:rounding}, componentwise truncation $Q_q$ toward zero to $\Cg_q$ satisfies
\begin{equation}\label{eq:walsh-truncation}
  \|Q_qz\|_\infty\le\|z\|_\infty,\qquad \|Q_qz-z\|_\infty<\sqrt2\cdot2^{-q}.
\end{equation}

We invoke the simultaneous butterfly layer through Proposition~\ref{prop:prop14-param}, proved in
Appendix~\ref{app:parameter-transfer} from the elementary stream primitives of Fact~\ref{fact:oai-streams} and the
fixed triple networks at $h=100$, $m=h^3=10^6$ of Fact~\ref{fact:oai-networks}. Let us recall its signature at grid
precision~$q$. Fix rational network exponents $0<\tau<1$ and $0\le\sigma<1$ satisfying $s_b/W_b<m^\tau$ and
$s_c/W_c<m^\sigma$, rational parameters $0<\beta<1$, $c>0$, $0<\eps<1$, $\lambda$, $\lambda'$ with
$\max\{\tau,\sigma,\tau(1+c/\beta)\}<\lambda<1$ and $\max\{\lambda,\sigma+\beta(1-\sigma)\}<\lambda'<1$, a guard
exponent $C_1\ge14$ with $\eps C_1<1$, and band constants $a_d\le b_d$, $a_r\le b_r$, $C_M$,
$C_{\rm desc}$ and $C_w>1$. For all $q$ above a fixed cutoff with
\[
  a_dq^{\eps}\le d\le b_dq^{\eps},\qquad 2^{a_rq/d}\le r\le 2^{b_rq/d},\qquad K=\lfloor d^c\rfloor,
\]
and integers $1\le D\le d$, $A,S\ge1$ and $0\le\rho<K$ satisfying $M=AS\,2^{KD}$, $\lg(2M)\le C_Mq$ and
$q+2\le w\le C_wq$, the procedure accepts a disk-grid array in lexicographic order on $[A]\times[2^K]^D\times[S]$
together with the header $\Gamma(q)\Gamma(d)\Gamma(r)\Gamma(D)\Gamma(A)\Gamma(S)\Gamma(K)\Gamma(\rho+1)\Gamma(w)$ of
length at most $C_{\rm desc}q$. Each address carries a length-$r$ record in $\Rr_r$ stored in coefficient order, each
real or imaginary component being $2^{-q}$ times a signed $w$-bit two's-complement integer written most significant bit
first. Let $H$ apply $H_0$ simultaneously to bit~$\rho$ (counted from the least significant end) of each of the $D$
active chunks. Proposition~\ref{prop:prop14-param} evaluates $Hz$ \emph{exactly} over $\Z[i,1/2]$ using
$\Delta=\lceil C_0d^{C_1}\rceil=o(q)$ internal guard bits ($q+\Delta$ fractional and $\Delta+3$ integer and sign bits
per real component) before applying one final truncation $Q_q$. It returns $Q_q(Hz)$ in the original record and
component format, as a disk array satisfying $\|Q_q(Hz)-Hz\|_\infty<\sqrt2\cdot2^{-q}$, in $O(Vd^{\lambda'})$ steps,
where $V=2Mrw$.

\subsection{Selective Walsh operations}\label{sec:walsh-select}

Controlled flips require $H_0$ on an arbitrary subset of a \emph{slice}, the set of $D$ positions selected by one call
of the butterfly layer (Definition~\ref{def:walsh-positional} fixes the slices used for positional maps); three full
layers provide this selective operator exactly. Positions in a binary record address are numbered from the most
significant end, while within each $K$-bit chunk the offset $\rho$ is counted from the least significant end. Let
$S=\diag(1,i)$, and for any set $U$ of record-address positions define the recordwise diagonal
\begin{equation}\label{eq:walsh-SU}
  (S_Uz)(x)=i^{\sum_{j\in U}x_j}\,z(x),
\end{equation}
which multiplies every coefficient of the record at address~$x$ by the same fourth-root phase.

\begin{lemma}[Three-layer selection]\label{lem:walsh-select}
Let a slice consist of $D$ selected positions ($1\le D\le d$) and let $T$ be any subset of them, with $s=|T|$, $U$ the
selected positions of the slice not in~$T$, and $q_{\rm sel}=|U|=D-s$. Let $H_{\rm all}$ apply $H_0$ on all $D$
positions and $H_T$ only on those in~$T$. Then
\begin{equation}\label{eq:walsh-three-layer}
  H_0^3=\tfrac12H_0,\qquad (H_0S)^3=\tfrac{1+i}{4}I,\qquad
  H_T=2^s\,[2(1-i)]^{q_{\rm sel}}\,(H_{\rm all}S_U)^3 .
\end{equation}
If $q_{\rm sel}=2f+e$ with $e\in\{0,1\}$, then $2^s[2(1-i)]^{q_{\rm sel}}=2^{s+3f+e}(-i)^f(1-i)^e$.
\end{lemma}

\begin{proof}
The first identity follows from $H_0^2=\frac12I$, and cubing
$H_0S=\frac12\bigl(\begin{smallmatrix}1&i\\1&-i\end{smallmatrix}\bigr)$ gives the second. Operators on different
positions commute and tensor independently, so $(H_{\rm all}S_U)^3$ acts position by position: as $\frac12H_0$ on each
position in~$T$ and as $\frac{1+i}{4}I$ on each position in~$U$. Hence
$(H_{\rm all}S_U)^3=2^{-s}(\frac{1+i}{4})^{q_{\rm sel}}H_T$, and $(\frac{1+i}{4})^{-1}=2(1-i)$. Finally,
$[2(1-i)]^2=-8i$ gives the last formula.
\end{proof}

In chronological order, the \emph{gadget} of Lemma~\ref{lem:walsh-select} applies
$S_U,H_{\rm all},S_U,H_{\rm all},S_U,H_{\rm all}$ and then the scalar $2^s[2(1-i)]^{q_{\rm sel}}$; there are three
$S_U$ diagonals, including the initial one. The scalar is applied in one scan: a fourth-root phase, at most one
multiplication by $1-i$ via $(\Re z,\Im z)\mapsto(\Re z+\Im z,\Im z-\Re z)$, and one binary shift. Every external
diagonal and scalar correction is applied only at completed layer boundaries in the ordinary record layout, after the
layer call has restored its role routing and removed its padding, with position masks kept as caller data on separate
tapes.

\begin{lemma}[Exact boundaries]\label{lem:walsh-exact}
Let the input be a disk array on the $q$-grid and let a slice have $1\le D\le d$ selected positions. Put $q_*=q+8d$,
lift the representation exactly to precision $q_*$, and execute the gadget of Lemma~\ref{lem:walsh-select} by completed
calls of Proposition~\ref{prop:prop14-param} at precision~$q_*$, each with its internal guard~$\Delta$. Then the gadget
is exact: every final truncation $Q_{q_*}$ of a layer call discards only zeros, and the output is $H_Tz$, on a grid of
denominator exponent at most $q+s$. After any diagonal with entries in $\{1,-1,i,-i\}$ (which preserves both the grid
and the disk), a second gadget with the same $T$ is again exact, with output on a grid of denominator exponent at most
$q+2s$. All completed outputs are disk arrays.
\end{lemma}

\begin{proof}
A full normalized $D$-bit layer has Gaussian dyadic entries with denominator dividing~$2^D$, and $S_U$ adds no
denominator bits. The three layer outputs of the first gadget therefore lie on grids of denominator exponent at most
$q+D$, $q+2D$ and $q+3D$, all strictly below $q_*=q+8d$, so each $Q_{q_*}$ is the identity; the internal computation of
each layer is exact by Proposition~\ref{prop:prop14-param}, whose internal denominators are controlled by~$\Delta$.
After the exact scalar scan, the value equals $H_Tz$, which has denominator exponent at most $q+s$. A diagonal with
entries in $\{1,-1,i,-i\}$ multiplies every coefficient by a unit of $\Z[i]$, preserving both this denominator exponent
and the modulus of every coefficient. The second gadget therefore again starts from a disk array, now on denominator
exponent at most $q+s$; its three layer outputs have denominator exponents at most $q+s+3D\le q+4D<q_*$, and its
completed output $H_T(\cdot)$ has exponent at most $q+2s$. Each full layer and each $S_U$ is a contraction in
$\|\cdot\|_\infty$, and the completed gadget output equals $H_T$ applied to a disk array, hence is itself a disk array:
this operator identity, rather than a bound on the intermediate scalar, verifies the disk-input condition of the next
layer call. The scalar scan requires only a constant number of temporary integer bits.
\end{proof}

\begin{lemma}[Lifting]\label{lem:walsh-lift}
Lifting a $q$-grid array to precision $q_*=q+8d$ costs $O(V)$: append $8d$ fractional zero bits to every numerator and
use component width $w_*=w+8d$, so that $q_*+2\le w_*\le C_wq_*$, and rebuild the header with $q_*$, $w_*$ and the same
$d$, $r$, $K$. Since $d=O(q^\eps)$ with $\eps<1$, eventually $q_*\le2q$, and the comparison bands
$a_dq^\eps\le d\le b_dq^\eps$ and $2^{a_rq/d}\le r\le 2^{b_rq/d}$ at~$q$ imply
$(a_d/2)q_*^\eps\le d\le b_dq_*^\eps$ and $2^{(a_r/2)q_*/d}\le r\le2^{b_rq_*/d}$. Hence one fixed set of widened
comparison constants serves both $q$ and $q_*$.
\end{lemma}

\begin{proof}
Because $C_w>1$ and $w\le C_wq$, we have $w_*=w+8d\le C_w(q+8d)=C_wq_*$. Appending $8d$ low zero bits represents the
exact same dyadic values at scale $2^{-q_*}$ (whereas changing only the precision field in the header would rescale the
array), and the internal guarded width $q_*+2\Delta+3$ remains $O(q)$. All re-encoding scans are linear in~$V$.
\end{proof}

\subsection{Exact controlled flips}\label{sec:walsh-cnot}

Two selective layers around a sign diagonal perform a batch of controlled flips exactly. For disjoint pairs
$(c_j,t_j)$, $1\le j\le s$, of distinct record-address positions, let $T=\{t_1,\dots,t_s\}$ and define
\begin{equation}\label{eq:walsh-CZ}
  (\CZ\,z)(x)=(-1)^{\sum_{j=1}^sx_{c_j}x_{t_j}}\,z(x).
\end{equation}

\begin{lemma}[Controlled flips]\label{lem:cnot}
As array operators, $\CNOT^{\otimes s}=2^sH_T\,\CZ\,H_T$, where the left-hand side moves the record at address $x$ to
the address obtained by $x_{t_j}\mapsto x_{t_j}\oplus x_{c_j}$ for all $1\le j\le s$ and changes no coefficient.
\end{lemma}

\begin{proof}
For a single pair $(c_1,t_1)$, fix the control bit $x_{c_1}$. If $x_{c_1}=0$, the target diagonal is $I$ and
$2H_0^2=I$. If $x_{c_1}=1$, the target diagonal is $\diag(1,-1)$ and
$2H_0\diag(1,-1)H_0=\bigl(\begin{smallmatrix}0&1\\1&0\end{smallmatrix}\bigr)$. Tensoring over the $s$ disjoint pairs
gives the identity. Because the address map is an involution, the move and pull conventions coincide.
\end{proof}

\begin{proposition}[Slice-batched controlled flips]\label{prop:cnot-batch}
Assume the hypotheses of Proposition~\ref{prop:prop14-param} at grid precision~$q$. Suppose all targets $t_1,\dots,t_s$
of a batch of disjoint controlled flips lie among the selected positions of one slice, while the controls $c_1,\dots,c_s$
lie anywhere else in the record address (including selected positions of the same slice outside~$T$). On disk arrays on
the $q$-grid, the batch is performed exactly in $O(Vd^{\lambda'})$ steps by six completed layer calls at precision
$q_*=q+8d$ and a constant number of linear scans, returning the output in the original $q$-grid format.
\end{proposition}

\begin{proof}
Lift the input from $q$ to $q_*=q+8d$ (Lemma~\ref{lem:walsh-lift}), apply the first selective gadget for $H_T$, the sign
diagonal $\CZ$, and the second selective gadget for $H_T$, multiply exactly by $2^s$, and lower from $q_*$ back to~$q$.
Since $\CZ$ has entries in $\{\pm1\}$, Lemma~\ref{lem:walsh-exact} applies: all six layer truncations discard only
zeros. If a control $c_j$ lies in the same slice outside~$T$, the internal layers of a gadget mix that bit temporarily,
but the completed gadget acts on it as the identity before $\CZ$ is applied. By Lemma~\ref{lem:cnot}, the value after
multiplying by $2^s$ is an exact permutation of the input $q$-grid records, so lowering from $q_*$ to $q$ discards only
zero low bits; this relies on the permutation identity rather than denominator counting alone. There is no
control-before-target ordering requirement because $\CZ$ is a recordwise value diagonal rather than a controlled address
rotation. For each record, the phases of~\eqref{eq:walsh-SU} and the parity of~\eqref{eq:walsh-CZ} are computed once
from its address and the caller's mask or pair list in time polynomial in~$q$, which is absorbed by the record length
$r$ (Lemma~\ref{lem:transfer-long-record}), while lifting and lowering are elementary streams
(Lemma~\ref{lem:transfer-streams}) and the lifted volume remains $O(V)$. Thus Proposition~\ref{prop:prop14-param} gives
total time $O(Vd^{\lambda'})$.
\end{proof}

Because a completed batch is an exact record permutation on $\Cg_q$, it neither amplifies inherited numerical error nor
accumulates denominator bits across batches.

\subsection{Positional permutations}\label{sec:walsh-perm}

\begin{definition}[Positional maps and slices]\label{def:walsh-positional}
A \emph{positional map} on $m_0$ record-address positions is a permutation $\pi$ of $\{0,\dots,m_0-1\}$ under the move
convention: the bit at position $j$ is placed at position $\pi(j)$. It induces an array operator that permutes whole
records while leaving each record's coefficient list unchanged; $f\circ g$ applies $g$ first. Put
$J=\lfloor m_0/K\rfloor$, partition the first $JK$ positions into consecutive $K$-bit chunks, and group those chunks
into consecutive groups of at most $d$ chunks. For a group starting at position $j_0$ with $D$ chunks ($1\le D\le d$)
and an offset $0\le\rho<K$, the corresponding \emph{slice} consists of the positions $j_0+jK+(K-1-\rho)$ for
$0\le j<D$, with complete rectangular layer shape $[2^{j_0}]\times[2^K]^D\times[2^{m_0-j_0-DK}]$. The fewer than $K$
remaining address positions are \emph{leftover} positions.
\end{definition}

The number of slices is
\begin{equation}\label{eq:walsh-slices}
  N_{\rm slice}=K\lceil J/d\rceil\le \frac{m_0}{d}+K .
\end{equation}

\begin{lemma}[Two involutions]\label{lem:walsh-involutions}
Every positional map is a composition $I_1\circ I_0$ of two involutions. Each involution is a set of disjoint
positional swaps, and each set of disjoint swaps is implemented by three stages of disjoint controlled flips.
\end{lemma}

\begin{proof}
On a cycle $(i_0,\dots,i_{n_c-1})$ with $\pi(i_j)=i_{j+1}$ (indices modulo~$n_c$), define $I_0(i_j)=i_{-j}$ and
$I_1(i_j)=i_{1-j}$. Both are involutions and $I_1(I_0(i_j))=i_{j+1}$; defining them independently on every cycle of
$\pi$ gives $\pi=I_1\circ I_0$. A swap of two positions $b_1,b_2$ factors into three controlled flips,
$x_{b_1}\leftarrow x_{b_1}\oplus x_{b_2}$, $x_{b_2}\leftarrow x_{b_2}\oplus x_{b_1}$,
$x_{b_1}\leftarrow x_{b_1}\oplus x_{b_2}$; within each of the three stages, distinct swaps act on disjoint pairs of
positions and therefore commute.
\end{proof}

\begin{theorem}[Positional permutations]\label{thm:positional-perm}
Assume the parameter family of Proposition~\ref{prop:prop14-param} at grid precision~$q$. Let a complete $m_0$-bit
record-address cube, with $m_0=O(q)$, carry disk-valued $q$-grid records of $r$ complex coefficients, so that
$V=2^{m_0}\cdot2rw$. For every positional map~$\pi$ on $\{0,\dots,m_0-1\}$, a fixed multitape procedure performs the
induced record permutation exactly in
\begin{equation}\label{eq:walsh-perm-cost}
  O\Bigl(V\Bigl[\Bigl(\frac{m_0}{d}+K\Bigr)d^{\lambda'}+K\Bigr]\Bigr)
\end{equation}
steps, which is $O(Vqd^{\lambda'-1})$ when $dK=o(q)$. The output is returned in the original $q$-grid format.
\end{theorem}

\begin{proof}
\emph{Schedule.} Execute the three stages of $I_0$ followed by the three stages of $I_1$ from
Lemma~\ref{lem:walsh-involutions}. Within each stage, partition the controlled-flip pairs according to the slice that
contains their target bit; reordering batches within a stage is valid because the pairs of a stage have pairwise
disjoint supports. Each slice batch satisfies Proposition~\ref{prop:cnot-batch} (including pairs whose control bit lies
in a leftover position or in another slice). Any pair whose target bit is a leftover position is executed directly by a
two-digit stream permutation (Lemma~\ref{lem:transfer-streams}(1)); since there are fewer than $K$ leftover positions,
each of the six stages performs fewer than $K$ such permutations, costing $O(VK)$ in total.

\emph{Bookkeeping.} Cycle decompositions, involution tables, slice masks and nine-field headers have size polynomial
in~$q$ and are constructed by polynomially many scans of polynomial-size tables, absorbed into $O(V)$ by the record
length (Lemma~\ref{lem:transfer-long-record}). Every slice acts on the complete cube of $2^{m_0}$ records, its two
spectator ranges contribute at most $2m_0+2$ header bits, and the record length contributes $2\lg r+1=O(q)$ bits, so a
single pair of constants $C_M,C_{\rm desc}$ serves every slice call.

\emph{Cost.} Six stages with at most $N_{\rm slice}$ slice batches per stage and six layer calls per nonempty batch use
at most $36N_{\rm slice}$ completed layer calls. Combining~\eqref{eq:walsh-slices} with Proposition~\ref{prop:cnot-batch}
yields~\eqref{eq:walsh-perm-cost}. Each layer completes its own padding and cleanup, and each batch returns to the
$q$-grid before the next begins; internal packed basis changes (Lemmas~\ref{lem:transfer12}--\ref{lem:transfer13}) are
already accounted for in $O(Vd^{\lambda'})$. When $m_0=O(q)$ and $dK=o(q)$, we have
$(m_0/d+K)d^{\lambda'}+K=O(qd^{\lambda'-1})$.
\end{proof}

Note that Theorem~\ref{thm:positional-perm} permutes the bit positions of a complete binary record address carrying a
long contiguous payload via intermediate linear combinations on $\Cg_{q_*}$. It does not execute arbitrary permutations
of the $2^{m_0}$ records or apply to short payloads without a record factor~$r$.
